\section{What the framework repairs, and what it leaves open}
\label{sec:audit}

\subsection{Repairs supported by an explicit new argument}

\paragraph{The FTA iteration restriction.}
The baseline gives a one-root procedure uniform in degree, but restricts
its deflation paragraph to an externally fixed degree
\cite{repo-fta}. Theorem~\ref{thm:fta-uniform-factorization} supplies the
missing uniform nested-evaluation argument. It has more content than a
renaming of the schema: the output root evaluator has the numerical
degree as an argument, its stage caps are computed in advance, and its
verification horizons are composed before the stage induction.

The zeta-zero report already constructs a different rational-polynomial
root backend after identifying that restriction \cite{repo-zeta-roots}.
Our result extends to the fixed Cauchy-coefficient schema. It does not
replace the rational backend's exact gcd, multiplicity, or separation
computations; those require different finite data.

\paragraph{The Bohr--Mollerup candidate modulus.}
The source asks whether its extra candidate modulus follows from the
other hypotheses \cite{repo-bm}. Theorem~\ref{thm:convex-local-modulus}
answers this for the source's rational presentation. It uses strong
positivity and multiplicative convexity, and in fact does not need the
normalization or recurrence. The modulus is a concrete expression in
rational endpoints and two coarse evaluations. The finite AM--GM padding
construction also answers the report's request for a simpler arithmetic
certificate architecture; we make no claim of a measured reduction in
checker bytes or native proof length.

\paragraph{The finite geometric section bound.}
The Siegel route separates an approximation/cover statement from a
zero-count or fibre bound \cite{repo-siegel}. Theorem~\ref{thm:superelliptic-section}
supplies that bound for $y^r=f(x)$, $\deg f>rs$, with explicit features.
For the quintic family in the Faltings report \cite{repo-faltings}, the
bound is five. A valid cover would yield $5M$ as a cardinality bound.
The existence of such covers for the relevant set of points is not
proved here. For some unrestricted point sets it is false; the geometric
bound must therefore not be presented as a theorem of general finiteness.

\paragraph{The common analytic interfaces.}
The calculus and Taylor--Cauchy reports already state fixed rational
samplers and modulus hypotheses \cite{repo-calculus,repo-taylor}.
The finite-net theorem makes their useful compact-domain extremum module
available once, with a precise output-value name. The coherent-minimizer
theorem separates the additional quantitative uniqueness information
needed to obtain a point. This is an extension of shared infrastructure,
not a claim that those source arguments were incorrect.

\subsection{Obstructions that cannot be repaired by changing the encoding alone}

\paragraph{Exact Puiseux data from arbitrary coefficient names.}
The Puiseux report gives a halting obstruction already in the family
$Y^2-X^2-a$, with $a$ supplied only by a Cauchy name
\cite{repo-puiseux}. Even the promise that the value happens to be a
dyadic rational does not provide its exact finite presentation. A zero
versus positive coefficient can alter the exact initial branch data.
The finite-observation framework retains exact algebraic presentations
or explicit discrete certificates when the requested output requires
them. Returning coherent approximate complex roots, as our factorization
theorem does, does not decide those branch data.

\paragraph{Apartness and inversion.}
A procedure dividing by a represented number needs quantitative
nonvanishing. The source FTA discussion uses $az-1$ to explain this
boundary \cite{repo-fta}. Proposition~\ref{prop:name-arithmetic} supplies
an explicit reciprocal construction from an apartness certificate,
while denominator-free identities and monic presentations can avoid some
uses of division. No universal nonzero-to-apartness principle is added.

\paragraph{Uniform evaluation of all PR program codes.}
The diagonal argument in Section~\ref{sec:compiler} rules out a primitive
recursive evaluator for every primitive recursive code. The
zeta-functional-equation plan already recognizes this issue
\cite{repo-zeta-plan}. Our theorem solves a variable-depth composition
problem for a fixed uniform table routine; it does not invalidate that
obstruction. The difference between a numerical degree and an arbitrary
program code is essential.

\paragraph{Metastability, cardinality, and cutoffs.}
The one-jump sequence in Section~\ref{sec:quantitative-analysis} has a
uniform finite-window procedure but no uniform computable Cauchy rate.
The singleton predicate in Section~\ref{sec:finite-sections} has a
uniform cardinality bound and vacuous pairwise gaps but no computable
uniform height cutoff. These prove failures of general conversion
principles. They do not establish that any named Diophantine theorem is
unprovable in \iPRA, nor that a specific sequence cannot have a modulus.

\subsection{Analytic number theory and probability: exposed dependencies}

The following are examples of mathematical input still needed after the
finite compiler is available. They are not an exhaustive classification
of the proof-theoretic strength of the corresponding classical theorems.

\begin{longtable}{@{}>{\raggedright\arraybackslash}p{0.25\linewidth}>{\raggedright\arraybackslash}p{0.43\linewidth}>{\raggedright\arraybackslash}p{0.26\linewidth}@{}}
\toprule
Reported difficulty & What a shared finite framework supplies & What it does not supply\\
\midrule
\endhead
Stronger prime number theorem error & Rational contour/error bookkeeping,
precision allocation, and common finite induction invariants. & A
stronger zero-free-region estimate. The inspected report establishes its
own weaker exponent, not the optimal requested rate \cite{repo-pnt}.\\[5pt]
Siegel--Walfisz and Bombieri--Vinogradov uniformity & Explicit contracts
for exceptional terms and their parameter dependence. & An effective
Siegel-type separation estimate, or the unproved exceptional-zero formula
itself \cite{repo-prime-sw}.\\[5pt]
Chebotarev estimates & Finite field-data and analytic-certificate
interfaces. & The required field-dependent analytic estimates and
exceptional-zero control \cite{repo-chebotarev}.\\[5pt]
General-law and advanced central limit statements & Finite-law
calculations can transfer through supplied PR tail, quantization, or
Lindeberg moduli. & High-dimensional convex smoothing, entropy-power
arguments, or a full stable-law classification \cite{repo-clt}.\\
\bottomrule
\end{longtable}

One particularly important uniformity issue is that excluding a single
exceptional modulus does not exclude all moduli divisible by its
conductor. The inspected prime-distribution report explicitly keeps this
problem visible \cite{repo-prime-sw}. Our framework preserves the
quantifiers and parameter dependence; it does not simplify the claim by
dropping the exceptional multiples.

\subsection{A concrete integration sequence}

A useful repository integration would proceed in three pieces. First
add the finite data, Boolean reflection, and table-compilation modules
under a shared arithmetic namespace; keep them independent of a single
analytic theorem. Second add the factorization, convex-modulus, and
superelliptic-section applications with their source-specific hypotheses.
Third refactor existing packets to call those modules, preserving any
remaining mathematical contracts.

The supplied archive includes an integration note, not a repository
mutation. Its priority is a reviewable mathematical and formalization
package. Nothing in this report asserts that the existing repository's
uncompiled outlines have thereby become completed object-theory proofs.
